\renewcommand{\Anno}{2001}

\renewcommand{\NumeroProblema}{3}

\renewcommand{\Testo}{Comunque si scelgano quattro studenti della
Scuola S. Anna, ce n'è almeno uno dei quattro che è amico degli altri tre.  Dimostrare che esiste almeno uno studente della Scuola che è amico di tutti gli altri.
	
{\small Nota.  Si usi la convenzione che $A$ amico di $B$ implica $B$
amico di $A$.}}

\renewcommand{\Risposta}{\ldots}

\renewcommand{\Aiutino}{Conviene procedere per casi.  Supponiamo che
ci sia un $A$ che non conosce un certo $B$.  Prendiamo ora un terzo studente $C$.  Abbiamo 3 casi: o $C$ conosce tutti e allora abbiamo finito, o $C$ non conosce qualcuno diverso da $A$ e $B$ (ma in questo caso \ldots), o $C$ non conosce almeno uno tra $A$ e $B$, diciamo $A$.
Prendiamo ora un quarto studente $D$: chi può non conoscere $D$?}

\renewcommand{\SoluzioneUfficiale}{
}

\renewcommand{\SoluzioniAlternative}{
}

\renewcommand{\Osservazioni}{
}

\renewcommand{\Commenti}{
}